\documentclass[11pt, a4paper]{article}
\usepackage[utf8]{inputenc}
\usepackage[spanish,es-noshorthands]{babel}
\usepackage[margin=2cm]{geometry}
\usepackage{enumitem}
\usepackage{titlesec}
\usepackage{xcolor}

\definecolor{primary}{RGB}{0, 80, 160}

\titleformat{\section}
  {\Large\bfseries\color{primary}}
  {}{0em}
  {}
  [\titlerule]

\begin{document}

\begin{center}
    {\huge \textbf{César Augusto Rodríguez Gómez}}\\[0.2cm]
    {\large Investigador en Electrónica e Inteligencia Artificial Multimodal}\\[0.2cm]
    \texttt{cero97@yahoo.com} $\vert$ \texttt{+58 0412-188.31.84} $\vert$ \texttt{github.com/cero97-ctrl} $\vert$ \texttt{Telegram: @cero2k6}
\end{center}

\vspace{0.5cm}

\section*{Perfil Profesional}
Lic en Física/Investigador especializado en la convergencia de la Ingeniería Electrónica, el Internet de las Cosas (IoT) y la Inteligencia Artificial. Con sólida experiencia en la creación de flujos de trabajo automatizados, diseño de arquitecturas de agentes LLM en tres capas y sistemas de Generación Aumentada por Recuperación (RAG) aplicados a la Automatización de Diseño Electrónico (EDA) y al sector educativo. Miembro y coordinador del Grupo de Investigación GIDEAL en la Universidad de Oriente.

\section*{Experiencia en Investigación y Desarrollo}

\noindent \textbf{Investigador Principal / Desarrollador de Inteligencia Artificial} \hfill \textit{Actualidad} \\
\textit{Grupo de Investigación GIDEAL -- Universidad de Oriente}
\begin{itemize}[leftmargin=1.5em, itemsep=0.3em]
    \item \textbf{Orquestador Multimodal de Evaluación:} Desarrollo de un sistema basado en IA (Gemini y OpenRouter) para analizar fotografías de circuitos y hacer evaluaciones preliminares (sujetas a revisión) de exámenes manuscritos y reportes de laboratorio.
    \item \textbf{Agente Inteligente de Diseño Electrónico (EDA):} Creación de una herramienta algorítmica de ingeniería inversa capaz de interpretar código \LaTeX{} (paquete \texttt{circuitikz}) para la extracción automática de \textit{netlists} y Listas de Materiales (BOM) exportables a KiCAD y EasyEDA.
    \item \textbf{Asistente Académico (Sistemas RAG):} Implementación de bases de datos vectoriales (ChromaDB) y grandes modelos de lenguaje (Groq) para el procesamiento conversacional y veloz de extensa documentación técnica e instruccional en PDF y Markdown.
    \item \textbf{Automatización y Orquestación Remota:} Diseño de una arquitectura distribuida de servidores bajo el estándar MCP (Model Context Protocol), integrando una pasarela (Gateway) segura vía Telegram para acceder a flujos experimentales sin redes complejas.
    \item \textbf{Automatización de Documentación Científica:} Orquestación de scripts para la emisión automática de problemarios, rúbricas y artículos compilables directamente en \LaTeX.
\end{itemize}

\section*{Habilidades y Herramientas Técnicas}
\begin{itemize}[leftmargin=1.5em, itemsep=0.3em]
    \item \textbf{Desarrollo de Software y SO:} Python (Avanzado), \LaTeX, Bash Scripting, Linux, Systemd, Git.
    \item \textbf{Inteligencia Artificial y NLP (Procesamiento de Lenguaje Natural):} Arquitectura de Agentes de 3 Capas, Grandes Modelos de Lenguaje (LLMs de código abierto y cerrados), LangChain, ChromaDB, Prompt Engineering Multimodal.
    \item \textbf{Ingeniería Electrónica:} Diseño de Circuitos Impresos (PCB), Internet de las Cosas (IoT), Fuentes Conmutadas (SMPS), Software EDA (EasyEDA, KiCAD, simuladores).
\end{itemize}

\section*{Educación}
\noindent \textbf{Título Universitario (Lic. en Física / Electrónica)} \hfill \textit{1991} \\
\textit{Universidad de Oriente, Núcleo Sucre, Venezuela}

\end{document}
